\documentclass{article}
\usepackage{amsmath,amssymb,amsthm}
\newtheorem{theorem}{Theorem}
\begin{document}
\title{Polynomial-size shared correction of moment-curve jets}
\author{}
\date{November 24, 2026}
\maketitle

\paragraph{Verification status.}
The arbitrary independent-column polynomial correction theorem, its exact
head count, numerical limiting identity, full-cube denominator positivity,
strictly negative slopes, zero-column axis, and preservation of strict
signs are Lean checked in \texttt{JointFamilyRidgeCorrections.lean}.
The moment-curve interpolation specialization and feature-space analysis
below are complete written arguments. The finite-field rank data are
computations over the specified finite field, not Lean proofs or assertions
of equal real rank. No improved global upper bound is claimed.

\section{One polynomial costs one more head than its degree}
Let $L(x)=\sum_{i=1}^k w_i x_i$, where every $w_i>0$. Each column
$j$ chooses an arbitrary real polynomial $Q_j$ of degree at most $R$.
Put $H=R+1$ and choose distinct exterior poles
$p_h=T+h$, for $h=0,\ldots,R$. Barycentric partial fractions give
real residues $c_{hj}$ satisfying the exact identity
\[
 \sum_{h=0}^R\frac{c_{hj}}{p_h-L(x)}
 =Q_j(L(x))\prod_{h=0}^R\frac{p_h}{p_h-L(x)}.
 \tag{1}
\]
For example, with $N(t)=\prod_h(t-p_h)$ one may take
\[
 c_{hj}=-N(0)Q_j(p_h)\prod_{\ell\ne h}(p_h-p_\ell)^{-1}.
\]
The product in (1) tends to one as $T$ tends to infinity, uniformly
on the finite row cube. It is the same product for every column.

Use the numerator $\sum_jc_{hj}z_j$ and denominator
\[
 T+h+1-\sum_iw_ix_i-\sum_jz_j.
\]
When $T>\sum_iw_i+s$, every denominator is positive on the entire
ambient cube, and every slope is strictly negative. Numerators are affine
and vanish on the zero-column axis. At a onehot column input the
denominator is exactly $p_h-L(x)$. The residues may be arbitrarily large
and may have either sign, as allowed by the original model.

\begin{theorem}
Given any residual score $G(x,j)$ for which
$G(x,j)+Q_j(L(x))$ has the required strict sign at every input, there
are $R+1$ admissible reverse heads whose score, added to $G$, preserves
all those strict signs simultaneously. For $m$ positive linear row
statistics and independently chosen degree-$R$ polynomials for each
statistic and column, the cost is exactly $m(R+1)$ heads.
\end{theorem}
\begin{proof}
Equation (1) converges to the polynomial correction at every input.
Finiteness and strict margins give a common sufficiently large $T$.
For several statistics, apply this argument successively, including
the remaining polynomial terms and all previously selected heads in
the residual. Concatenation preserves full-cube admissibility and the
exact zero-column axis. The main declarations are
\texttt{jointFamily\_polynomial\_ridge\_correction\_preserves\_signs}
and
\texttt{jointFamily\_polynomial\_ridge\_family\_correction\_preserves\_signs}.
\end{proof}

\section{Moment-curve specialization}
Put $P_x(t)=\sum_{i=1}^k x_it^i$. Let $V_R$ be the span, as
functions of $x$, of all coefficient functions of $P_x(t)^a$ for
$0\leq a\leq R$. Choose $kR+1$ distinct positive real numbers
$t_1,\ldots,t_{kR+1}$. Polynomial interpolation in $t$ shows that
every coefficient of $P_x(t)^a$ is a fixed linear combination of
$P_x(t_\ell)^a$, because its degree in $t$ is at most $kR$.
Consequently, every independently chosen column correction
$F_j\in V_R$ has a representation
\[
 F_j(x)=\sum_{\ell=1}^{kR+1}Q_{\ell j}(P_x(t_\ell)),
 \qquad \deg Q_{\ell j}\leq R.
\]
The row weights $t_\ell^i$ are strictly positive, so the theorem gives
the explicit shared correction cost
\[
 (kR+1)(R+1)=O(kR^2).
 \tag{2}
\]
No relation is imposed between corrections in different columns.
This permits subtracting every low-order jet of each permutation
separately, while retaining all $H!$ permutations. In contrast, restricting
permutations to a common exact integer-moment fiber pays a signature
entropy loss $O(k^2R^3\log H)$.

\section{Why this does not yet give a decoder}
For a permutation difference of reciprocal banks, the jet of order
$r$ is a linear combination of coefficient functions of $P_x(t)^a$
with $1\leq a\leq r-1$. The terms with $a=0$ cancel because the
memory histogram is fixed; those with $a=r$ cancel because row
directions do not change. Thus the cumulative jet span through order
$r$ has dimension at most
\[
 \sum_{a=1}^{r-1}(ka+1)=\frac{kr(r-1)}2+r-1.
 \tag{3}
\]
The newly available span at the next order adds at most $kr+1$
dimensions. Subtracting lower jets does not automatically make the
residual covariance isotropic on a large input space: after rescaling
the first surviving order, all higher-order directions still vanish
in the small-scale limit. Its full leading image has dimension at most
$O(kr^2)$, so spanning $M$ input directions already requires correction
cost of order $M$. If one additionally quotients by the entire lower-jet
feature space, the first new image has dimension at most $kr+1$.
Spanning an $M$-dimensional remaining quotient in that single order
would require an index of order $M/k$ and correction cost of order
$M^2/k$. These estimates defeat the indicated cheap one-order arguments
in the exponential-domain regime. They are not lower bounds for
permutation decoders using several target-dependent cancellation orders.

There is a precise test for such multiple orders. If a vector $z$ is
nonzero on a set of required inputs and annihilates every jet through
order $R$ there, any reciprocal-difference germ having strict signs
$\operatorname{sign}(z)$ on that set for all sufficiently small positive
scales must have each of those jets vanish pointwise on that set.
Indeed, otherwise take the first nonzero order among the finitely many
coordinates. All its signed leading coefficients are nonnegative and
at least one is positive, contradicting their zero $z$-weighted sum.
This forces deep cancellations for appropriate targets; it does not
prove that such cancellations are impossible.

\paragraph{Finite evidence.}
The script \texttt{JointFamilyMomentJetRanks.py} computes the exact
ranks of the spans generated by $P_x(t)^a$, $1\leq a\leq R$, on
all nonzero Boolean inputs, over the prime field of order $65521$.
For $k=4,5,6,7,8,9,10$, full finite-field rank first occurred at
$R=4,7,8,10,11,14,18$, respectively. For example, at $k=10,R=10$
the rank is $387$ on $1023$ inputs, and at $R=17$ it is $981$.
These computations illustrate the distinct surviving orders, but
do not establish an asymptotic estimate or exact real rank.

\end{document}
